\documentclass[12pt,a4paper]{article}
\usepackage[T1]{fontenc}
\usepackage[utf8]{inputenc}
\usepackage[polish]{babel}
\usepackage{lmodern}
\usepackage{geometry}
\usepackage{setspace}
\usepackage{amsmath,amssymb}
\usepackage{graphicx}
\usepackage{float}
\usepackage{booktabs}
\usepackage{caption}
\usepackage{listings}
\usepackage{xcolor}
\usepackage{hyperref}
\usepackage{array}
\usepackage{multirow}

\geometry{margin=2.5cm}
\onehalfspacing
\captionsetup{font=small,labelfont=bf}

\lstdefinestyle{pythonstyle}{
    language=Python,
    basicstyle=\ttfamily\small,
    keywordstyle=\color{blue},
    commentstyle=\color{teal},
    stringstyle=\color{purple},
    showstringspaces=false,
    frame=single,
    breaklines=true
}

\begin{document}

\begin{titlepage}
    \centering
    \includegraphics[width=0.35\linewidth]{logo_uczelni.png}\par
    \vspace{1.0cm}
    {\large Politechnika / Uczelnia \par}
    \vspace{1.5cm}
    {\Large\bfseries Wieloetapowe procesy decyzyjne, Laboratorium 07\par}
    \vspace{1.2cm}
    {\large Przedmiot: Wprowadzenie do Przetwarzania Danych / WPD\par}
    \vspace{2cm}
    \begin{flushleft}
        \textbf{Autor:} Stanis\l{}aw Krupa \\
        \textbf{Nr albumu:} 254977 \\
        \textbf{Grupa:} XX \\
        \textbf{Prowadz\k{a}cy:} dr in\.z. Marlena Dr\k{a}g \\
        \textbf{Data wykonania:} 15.06.2026 \\
    \end{flushleft}
    \vfill
    {\large \today \par}
\end{titlepage}

% ─────────────────────────────────────────────────────────────────
\section{Cel pracy}
Celem laboratorium by\l{}o wyznaczenie optymalnej strategii zakupu
$(u_1^*,\ldots,u_n^*)$ minimalizuj\k{a}cej \l{}\k{a}czny koszt zakupu surowca
w wielookresowym zagadnieniu zarz\k{a}dzania magazynem, zgodnie z~przyk\l{}adem
18.3.6.1 podr\k{e}cznika. Do~rozwi\k{a}zania zastosowano metod\k{e} r\'owna\'n
funkcjonalnych Bellmana (dyskretna optymalizacja dynamiczna).

% ─────────────────────────────────────────────────────────────────
\section{Opis zadania}

\subsection*{Dane problemu}
Rozpatrywano fab\-ryk\k{e} podzielon\k{a} na $n$ pr\.zed\-zia\l{}\'ow
odcinka czasowego. W~ka\.zdym okresie $j$ nale\.zy zakupi\'c $u_j$ jednostek
surowca po cenie $c_j$ za~jednostk\k{e}, tak aby zaspokoi\'c zapotrzebowanie
$v_j$ i~nie przekroczy\'c pojemno\'sci magazynu $K$.
Dane liczbowe przyk\l{}adu 18.3.6.1:

\begin{center}
\begin{tabular}{lcccccccc}
\toprule
Okres $j$ & -- & 1 & 2 & 3 & 4 & 5 & 6 \\
\midrule
Cena $c_j$ & & 4 & 3 & 5 & 3 & 4 & 2 \\
Zapotrzebowanie $v_j$ & & 6 & 7 & 4 & 2 & 4 & 3 \\
\bottomrule
\end{tabular}
\end{center}

\noindent
Pojemno\'s\'c magazynu $K=10$, pocz\k{a}tkowy stan magazynu $x_0 = x_a = 2$.

% ─────────────────────────────────────────────────────────────────
\section{Podstawy teoretyczne}

\subsection{Sformu\l{}owanie jako n-poziomowy proces decyzyjny}

Problem zakupu formalizujemy nast\k{e}puj\k{a}co:
\begin{equation}\label{eq:fc}
    \text{FC:}\quad f(u_1,\ldots,u_n) = \sum_{j=1}^{n} c_j u_j \;\longrightarrow\; \min,
\end{equation}
\begin{equation}\label{eq:wd}
    \text{WD:}\quad
    \begin{cases}
        x_j = x_{j-1} + u_j - v_j, & j = 1(1)n, \\
        x_0 = x_a,\quad 0 \le x_j \le K, & j = 1(1)n, \\
        U_j(x_{j-1}) = \bigl\{u_j : \max(0,\,v_j - x_{j-1}) \le u_j \le K - x_{j-1}\bigr\}, & j = 1(1)n.
    \end{cases}
\end{equation}

Interpretacja warunk\'ow: $u_j \ge v_j - x_{j-1}$ zapewnia zaspokojenie popytu
(stan magazynu nie jest ujemny), natomiast $u_j \le K - x_{j-1}$ gwarantuje
nieprzekroczenie pojemno\'sci $K$.

\subsection{R\'ownania funkcjonalne Bellmana}

Definiujemy funkcj\k{e} koszt\'ow cz\k{e}\'sciowych:
\begin{equation}\label{eq:bell_def}
    \phi_j(x_{j-1}) = \min_{\substack{u_j,\ldots,u_n \\ \text{dopuszczalne}}}
        \sum_{k=j}^{n} c_k u_k,
\end{equation}
czyli minimalny koszt od okresu $j$ do ko\'nca, przy stanie magazynu
$x_{j-1}$ przed okresem $j$.

Warunek ko\'ncowy i rekurencja Bellmana:
\begin{align}
    \phi_{n+1}(x_n) &= 0, \label{eq:bell_end}\\
    \phi_j(x_{j-1}) &= \min_{u_j \in U_j(x_{j-1})}
        \Bigl[ c_j u_j + \phi_{j+1}(x_{j-1} + u_j - v_j) \Bigr],
        \quad j = n, n{-}1, \ldots, 1. \label{eq:bell_rec}
\end{align}

Optymaln\k{a} strategi\k{e} odtwarza si\k{e} iteracj\k{a} w~prz\'od:
pocz\k{a}wszy od $x_0 = x_a$, w~ka\.zdym kroku $j$ wybieramy
$u_j^* = \arg\min$~wyra\.zenia~\eqref{eq:bell_rec} i~obliczamy
$x_j = x_{j-1} + u_j^* - v_j$.

% ─────────────────────────────────────────────────────────────────
\section{Przebieg oblicze\'n}

\subsection{Iteracja wstecz -- wyznaczanie $\phi_j$}

\paragraph{Krok $j = 6$ \quad ($c_6 = 2$, $v_6 = 3$, $\phi_7 \equiv 0$).}
\[
    \phi_6(x_5) = \min_{u_6 \in U_6(x_5)} 2\,u_6
    = 2 \cdot \max(0,\, 3 - x_5).
\]
Wyb\'or: $u_6^* = \max(0,\,3-x_5)$ -- kupujemy dok\l{}adnie brakuj\k{a}ce
do zaspokojenia popytu, bo dalszych okres\'ow nie ma.
Warto\'sci: $\phi_6(0)=6$,\; $\phi_6(1)=4$,\; $\phi_6(2)=2$,\;
$\phi_6(x_5)=0$ dla $x_5 \ge 3$.

\paragraph{Krok $j = 5$ \quad ($c_5 = 4$, $v_5 = 4$).}
\[
    \phi_5(x_4) = \min_{u_5 \ge \max(0,4-x_4)}^{10-x_4}
        \bigl[4\,u_5 + \phi_6(x_4 + u_5 - 4)\bigr].
\]
Przyk\l{}adowo dla $x_4 = 0$: $u_5 \in [4,10]$;
najta\'nsze $u_5 = 4$: $4{\cdot}4 + \phi_6(0) = 16{+}6 = 22$.
Dla $x_4 = 4$: $u_5 = 0$ daje $0 + \phi_6(0) = 6$.

\paragraph{Krok $j = 4$ \quad ($c_4 = 3$, $v_4 = 2$).}
Dla $x_3 = 0$: sprawdzamy $u_4 \in [2, 10]$:
\begin{alignat*}{3}
    u_4=2{:}\ &3{\cdot}2 + \phi_5(0) = 6{+}22 = 28, &\qquad
    u_4=6{:}\ &3{\cdot}6 + \phi_5(4) = 18{+}6 = 24\ \leftarrow \min.
\end{alignat*}
Warto\'s\'c $\phi_4(0) = 24$, $u_4^* = 6$.
Zakup 6~jednostek (zamiast minimalnych~2) jest op\l{}acalny, gdy\.z
cena $c_4=3$ jest ni\.zsza ni\.z $c_5=4$ -- op\l{}aca si\k{e} zatowarowa\'c
przed dro\.zszym okresem~5.

\paragraph{Krok $j = 3$ \quad ($c_3 = 5$, $v_3 = 4$).}
Dla $x_2 = 0$: $u_3 \in [4,10]$;
$u_3 = 4$: $5{\cdot}4 + \phi_4(0) = 20{+}24 = 44\ \leftarrow \min$.
Cena $c_3 = 5$ jest najwy\.zsza -- nie op\l{}aca si\k{e} kupowa\'c z~nadwy\.zk\k{a}.

\paragraph{Krok $j = 2$ \quad ($c_2 = 3$, $v_2 = 7$).}
Dla $x_1 = 0$: $u_2 \in [7,10]$:
\begin{alignat*}{4}
    u_2=7{:}&\ 21+\phi_3(0)=65, &\ u_2=8{:}&\ 24+\phi_3(1)=63, \\
    u_2=9{:}&\ 27+\phi_3(2)=61, &\ u_2=10{:}&\ 30+\phi_3(3)=59\ \leftarrow \min.
\end{alignat*}
$\phi_2(0) = 59$, $u_2^* = 10$ -- maksymalne nape\l{}nienie magazynu,
poniewa\.z $c_2=3 < c_3=5$.

\paragraph{Krok $j = 1$ \quad ($c_1 = 4$, $v_1 = 6$).}
Dla $x_0 = 2$ (stan pocz\k{a}tkowy): $u_1 \in [4,8]$:
\begin{alignat*}{4}
    u_1=4{:}&\ 16+\phi_2(0)=75\ \leftarrow \min, &\ u_1=5{:}&\ 20+\phi_2(1)=76.
\end{alignat*}
$\phi_1(2) = 75$, $u_1^* = 4$.

\subsection{Tabele wynik\'ow iteracji wstecz}

Tabela~\ref{tab:phi} zawiera warto\'sci $\phi_j(x_{j-1})$ dla wszystkich
$j = 1,\ldots,6$ i $x_{j-1} = 0,\ldots,10$.
Tabela~\ref{tab:uopt} zawiera odpowiadaj\k{a}ce optymalne decyzje
$u_j^*(x_{j-1})$.

\begin{table}[H]
    \centering
    \caption{Warto\'sci funkcji Bellmana $\phi_j(x_{j-1})$}
    \label{tab:phi}
    \footnotesize
    \setlength{\tabcolsep}{5pt}
    \begin{tabular}{c|rrrrrrrrrrr}
        \toprule
        $j \backslash x$ & 0 & 1 & 2 & 3 & 4 & 5 & 6 & 7 & 8 & 9 & 10 \\
        \midrule
        1 & 83 & 79 & \textbf{75} & 71 & 67 & 63 & 59 & 56 & 53 & 50 & 47 \\
        2 & 59 & 56 & 53 & 50 & 47 & 44 & 41 & 38 & 35 & 32 & 29 \\
        3 & 44 & 39 & 34 & 29 & 24 & 21 & 18 & 15 & 12 &  9 &  6 \\
        4 & 24 & 21 & 18 & 15 & 12 &  9 &  6 &  4 &  2 &  0 &  0 \\
        5 & 22 & 18 & 14 & 10 &  6 &  4 &  2 &  0 &  0 &  0 &  0 \\
        6 &  6 &  4 &  2 &  0 &  0 &  0 &  0 &  0 &  0 &  0 &  0 \\
        \bottomrule
    \end{tabular}
\end{table}

\begin{table}[H]
    \centering
    \caption{Optymalne decyzje $u_j^*(x_{j-1})$}
    \label{tab:uopt}
    \footnotesize
    \setlength{\tabcolsep}{5pt}
    \begin{tabular}{c|rrrrrrrrrrr}
        \toprule
        $j \backslash x$ & 0 & 1 & 2 & 3 & 4 & 5 & 6 & 7 & 8 & 9 & 10 \\
        \midrule
        1 &  6 &  5 & \textbf{4} &  3 &  2 &  1 &  0 &  0 &  0 &  0 &  0 \\
        2 & 10 &  9 &  8 &  7 &  6 &  5 &  4 &  3 &  2 &  1 &  0 \\
        3 &  4 &  3 &  2 &  1 &  0 &  0 &  0 &  0 &  0 &  0 &  0 \\
        4 &  6 &  5 &  4 &  3 &  2 &  1 &  0 &  0 &  0 &  0 &  0 \\
        5 &  4 &  3 &  2 &  1 &  0 &  0 &  0 &  0 &  0 &  0 &  0 \\
        6 &  3 &  2 &  1 &  0 &  0 &  0 &  0 &  0 &  0 &  0 &  0 \\
        \bottomrule
    \end{tabular}
\end{table}

\subsection{Iteracja w prz\'od -- odtworzenie optymalnej strategii}

Startujemy od $x_0 = 2$ i korzystamy z~tabeli~\ref{tab:uopt}:

\begin{center}
\begin{tabular}{cccccc}
\toprule
Okres $j$ & $x_{j-1}$ & $u_j^*$ & $v_j$ & $x_j = x_{j-1}+u_j^*-v_j$ & Koszt cz. $c_j u_j^*$ \\
\midrule
1 & 2  & 4  & 6 & $2+4-6=0$   & $4\cdot4=16$  \\
2 & 0  & 10 & 7 & $0+10-7=3$  & $3\cdot10=30$ \\
3 & 3  & 1  & 4 & $3+1-4=0$   & $5\cdot1=5$   \\
4 & 0  & 6  & 2 & $0+6-2=4$   & $3\cdot6=18$  \\
5 & 4  & 0  & 4 & $4+0-4=0$   & $4\cdot0=0$   \\
6 & 0  & 3  & 3 & $0+3-3=0$   & $2\cdot3=6$   \\
\midrule
\multicolumn{5}{r}{\textbf{\l{}\k{a}czny koszt}} & \textbf{75} \\
\bottomrule
\end{tabular}
\end{center}

\bigskip
\noindent
\textbf{Optymalna strategia:}
\[
    \bigl(u_1^*,\, u_2^*,\, u_3^*,\, u_4^*,\, u_5^*,\, u_6^*\bigr)
    = (4,\; 10,\; 1,\; 6,\; 0,\; 3),
    \qquad \phi_1(2) = 75.
\]

% ─────────────────────────────────────────────────────────────────
\section{Wnioski}

\begin{enumerate}
    \item \textbf{Optymalny koszt} ca\l{}kowitego zakupu surowca wynosi
          $\phi_1(x_0)=75$ (jednostek pieni\k{e}\.znych).

    \item \textbf{Strategia zakupu} silnie wykorzystuje r\'o\.znice cen mi\k{e}dzy
          okresami. W~szczeg\'olno\'sci:
          \begin{itemize}
              \item w~okresie~2 ($c_2=3$, najta\'nszy obok~4 i~6) kupuje si\k{e}
                    maksymaln\k{a} dopuszczaln\k{a} ilo\'s\'c~($u_2^*=10$),
                    aby unikn\k{a}\'c drogich zakup\'ow w~okresie~3 ($c_3=5$);
              \item w~okresie~4 ($c_4=3$) kupuje si\k{e} z~nadwy\.zk\k{a}
                    ($u_4^*=6 > v_4=2$), by nie kupowa\'c w~dro\.zszym
                    okresie~5 ($c_5=4$) -- dzi\k{e}ki temu $u_5^*=0$;
              \item w~drogi\-ch okresach 1~i~3 ($c_1=4$, $c_3=5$) kupuje si\k{e}
                    dok\l{}adnie tyle, ile potrzeba b\k{a}d\'z nieznacznie wi\k{e}cej.
          \end{itemize}

    \item \textbf{Zasada optymalnych podstrategii Bellmana} zosta\l{}a
          potwierdzona w~praktyce: ka\.zda podstrategia $(u_j^*,\ldots,u_n^*)$
          jest r\'ownie\.z optymalna dla procesu cz\k{e}\'sciowego
          startuj\k{a}cego ze stanu $x_{j-1}$.

    \item Implementacja w~Pythonie (iteracja wstecz po macierzach
          \texttt{phi} i~\texttt{u\_opt}) daje identyczny wynik jak
          obliczenia r\k{e}czne, potwierdzaj\k{a}c poprawno\'s\'c obu podej\'s\'c.
\end{enumerate}

% ─────────────────────────────────────────────────────────────────
\section{Bibliografia}
\begin{enumerate}
    \item A.~Wr\'oblewski, \textit{Teoria optymalizacji -- Dyskretna
          optymalizacja dynamiczna}, rozdz.~18.3, wyd.~uczelni.
    \item R.~Bellman, \textit{Dynamic Programming}, Princeton University
          Press, 1957.
    \item Dokumentacja j\k{e}zyka Python 3, \url{https://docs.python.org/3/}.
\end{enumerate}

% ─────────────────────────────────────────────────────────────────
\section{Za\l{}\k{a}cznik -- kod \'zr\'od\l{}owy}

\begin{lstlisting}[style=pythonstyle, caption={Wyznaczanie optymalnej strategii zakupu -- przyk\l{}ad 18.3.6.1}]
n = 6
K = 10
x_a = 2

c = [None, 4, 3, 5, 3, 4, 2]   # ceny w okresach 1..6
v = [None, 6, 7, 4, 2, 4, 3]   # zapotrzebowania w okresach 1..6

phi   = [[float('inf')] * (K + 1) for _ in range(n + 2)]
u_opt = [[None]         * (K + 1) for _ in range(n + 2)]

for x in range(K + 1):
    phi[n + 1][x] = 0

# iteracja wstecz
for j in range(n, 0, -1):
    for x_prev in range(K + 1):
        u_min = max(0, v[j] - x_prev)
        u_max = K - x_prev
        for u in range(u_min, u_max + 1):
            x_j = x_prev + u - v[j]
            if 0 <= x_j <= K:
                cost = c[j] * u + phi[j + 1][x_j]
                if cost < phi[j][x_prev]:
                    phi[j][x_prev] = cost
                    u_opt[j][x_prev] = u

# odtwarzanie strategii (iteracja w przod)
policy, x_states, x_curr = [], [x_a], x_a
for j in range(1, n + 1):
    u_star = u_opt[j][x_curr]
    policy.append(u_star)
    x_curr = x_curr + u_star - v[j]
    x_states.append(x_curr)

print(f"Optymalny koszt: {phi[1][x_a]}")
print(f"Strategia u*:    {policy}")
print(f"Stany x:         {x_states}")
\end{lstlisting}

\end{document}
